\section{Evaluation}
The measurements below retain their archived executor and certificate results.
The subsequent three-state VC interface is checked by focused regressions;
archived unresolved mutations are not relabelled as new FAIL measurements.
\label{sec:evaluation}

The current evidence comprises executor regressions, a full QCP corpus scan,
an isolated deployment replay, and a first paired baseline smoke test. The
comparative study is governed by \texttt{research-protocol.md}; complete
Fulminate/CN experiments and proof-agent cost measurements remain in progress.

\subsection{Categorized executor regressions}

The development suite contains 52 functional programs with four inputs each,
three variants (reference, incorrect postcondition, and program mutation),
and two boundary policies. A further 37 diagnostic scenarios run under both
policies, yielding 1,322 checks across 18 categories. Functional expectations
are computed independently in Python and cross-checked against defined native
C executions. The engine itself uses symbolic addresses. The suite includes
779 negative checks, including intentionally wrong specifications and explicitly
unsupported execution scenarios.

All 1,322 recorded expectations are met: 541 runs return PASS, 757 UNKNOWN,
14 ERROR and 10 INVALID. No expected negative case is incorrectly accepted,
and no functional native oracle disagrees with the independent expected value.
Every one of the 52 body mutations is distinguished by at least one input.
UNKNOWN on an incorrect postcondition means that the engine retains an
obligation; it is not a proof of violation. The same matrix meets its
expectations under ASan/UBSan with leak detection disabled, and in a filesystem
namespace exposing only system libraries and the packaged teSpec tree.

The suite covers scalar control flow, C integer operations, loops, recursion,
calls, arrays, strings, aliases, structures, linked structures, trees and
sharing, allocator patterns, logical predicates, entry checks, safety,
declarations, incompleteness and annotation boundaries. It exposed an
existential-elimination bug that previously lost return constraints and caused
eight false acceptances. Preserving pre-existing witnesses during predicate
expansion removes those false acceptances; the associated negative cases now
retain return obligations.

\subsection{QCP corpus inventory}

A broad scan inventories 267 C files and 1,147 specifications and performs
5,160 executions, with another 287 cases requiring explicit typed bindings.
The execution outcomes are 1,086 PASS, 1,022 INVALID, 2,416 UNKNOWN and 636 ERROR.
These automatically generated small bindings are smoke inputs, not independently
labelled correctness tests. Entry invalidity and unresolved abstract predicates
therefore require separate interpretation from executor errors. Common execution
obstacles include unavailable readable/writable ownership, unbound proof
identifiers and annotation-branch handling. The comparison dataset will retain
these cases and report their feature categories and denominators.

\subsection{Baseline installation and first paired case}

CN is installed in a separate OCaml 5.2.0 environment at source revision
\texttt{181d0477}, with its pinned Cerberus dependency and Z3 4.13.3.
A smoke test uses the same C body in CN and QCP syntax: a function increments
an owned integer cell. The valid contract requires $0\leq v<100$ and returns
the cell with $v+1$; a mutated contract instead demands $v+2$.

\begin{center}
\small
\begin{tabular}{lll}
\toprule
Tool/task & Valid contract & Incorrect contract\\
\midrule
CN universal verification & Verified & Unprovable constraint\\
Fulminate, native cell $3$ & Check passes & Runtime check fails\\
\tespec{}, logical cell $3$ & PASS & UNKNOWN, residual VC\\
\bottomrule
\end{tabular}
\end{center}

The two instance tests agree on the initial cell value while \tespec{} leaves
the address symbolic. CN's verifier checks the broader precondition domain;
its result is recorded separately from matched instance checking. These six
runs validate installation and interfaces. They do not establish comparative
performance or overall superiority. Exact commands, outputs, source hashes and
resource measurements are retained by \texttt{run-comparison-smoke.py}.

Four further paired cases cover disjoint-cell exchange, structure fields,
nested calls and heap-updating loops. Each runs both postcondition variants
through the three tools; the call and loop cases additionally exercise
\tespec{}'s intermediate-boundary checking policy. These 36 additional runs
meet their installation expectations. The incorrect postconditions remain
unresolved in \tespec{} and fail Fulminate's runtime checks. These outcomes
motivate completing proof-backed failure classification before comparing
confirmed bug detection.

Two additional categories exercise shared aliases and pointer-valued storage.
In the first, two equal pointer arguments update the same owned integer; in
the second, the program loads a stored pointer and modifies its referent.
Identical C bodies and matched logical/native value 3 are used. Both \tespec{}
modes accept the valid contracts and leave the wrong postconditions UNKNOWN;
CN verifies the valid domains, and Fulminate detects the mutated postconditions
at runtime. Four unmodified-QCP runs discharge the valid targets and retain
one residual obligation per mutated target. The 24 CN/Fulminate/\tespec{}
invocations include repeated baseline runs for the two mode directories;
they represent two added subjects, not 24 independent subjects. Full logs and
VCs are archived in \texttt{memory-additions-20261011}. These concurrent probes
measure interface coverage; their resource records are not isolated timing
comparisons.

A recursive heap-update fixture adds a direct state-preservation probe. Its
helper increments an owned cell and recurses with a decreasing counter. With
logical value 3 and counter 3, \tespec{} enters the helper four times; boundary
mode also checks four helper contracts. Both modes pass with the exact helper
postcondition. Replacing that postcondition by an existential value bounded
below by the input preserves \tespec{}'s exact final value 6 and all selected
boundary checks. Fulminate also passes this concrete execution. Modular CN
verification and unmodified QCP retain an obligation for the top-level exact
postcondition because the helper summary is too weak. Thus the probe isolates
the effect of executing calls while checking contracts, rather than treating a
candidate contract as the next execution state. Mutating the top-level result
still produces a Fulminate runtime failure and a \tespec{} UNKNOWN result.
The 22 invocations and original weak-summary verification failures are retained
in \texttt{recursive-calls-20261011}; they cover two contract versions of one
C body.

\paragraph{Direct recursive-entry regression.}
Testing the recursive helper directly exposed an entry-specialization defect:
root bindings had also constrained its recursive call contract. A grid over
counters 0--3, values 0, 3 and $100-n$, three helper specifications and both
modes produced 18 unresolved valid boundary-mode cases without residual VCs.
Separating the named root instance from reusable contracts resolves all 48
valid-contract cases. The remaining 24 cases use a helper postcondition false
at the recursion base: boundary mode leaves all 12 unresolved with residual
obligations, while expand mode passes nine non-base cases with nested
annotations excluded and leaves the three base postconditions unresolved.
All 72 phase-aware expectations pass. The complete 1,322-case regression also
retains all expectations, including 779 negative checks with zero false passes.
The before/fixed outputs and the unsuccessful specification-ordering attempt
are retained in \texttt{entry-isolation-20261011}. These tests establish the
observed implementation behavior separately from the generic semantic proofs.

A further 25 instances exercise mutual recursion over a shared cell, including
an explicit selection of the original named contract on the return call. Both
functions update the cell before recursing. Valid cases preserve exact results
in both modes; a changed increment leaves obligations precisely on inputs that
reach the changed statement. Native outputs match independently calculated
values for both versions. These probes and a persistent integration test extend
the entry-isolation evidence to mutual and explicitly named calls.


An independently built, unmodified QCP executor at \texttt{e7745ada}, with
its pinned SMT revision \texttt{94d89c78}, checks the same five QCP-annotated
translation units over their full precondition domains. Four correct target
contracts are discharged automatically. The loop case retains a return
entailment whose invariant does not identify the current pointer with its
entry value. Adding a logical entry pointer and retaining that equality in
the invariant closes the remaining obligation; the C body is unchanged.
The original and refined annotations are both retained. All five incorrect
contracts, including the refined loop variant, retain proof obligations.
These are automatic solver results, separately labelled from kernel-checked
certificates: the upstream generated auto-proof files contain admitted
placeholders. The measurement covers all functions in the translation unit,
while outcomes are extracted specifically for the target function.

A separate April 2026 CN source snapshot at \texttt{39efc873} also runs the
paired cell checks successfully using its typed-integer annotation syntax.
Bennet and Darcy each complete ten generated tests of the valid contract at
seed zero and expose the incorrect contract with a reported cell containing
zero. Build dependencies and generated-harness compilers are recorded
separately. This source snapshot is distinguished from the archived paper
artifact, whose exact environment and original subjects are being reproduced.

The published v3 image has now been retrieved and checked against its published
digest. It contains the same CN revision and allocator revision
\texttt{2840aab7}, distinct from the later source-catalog snapshot. Using its
original user-space environment under rootless filesystem isolation on the
WSL2 host, the original quick script passes native Fulminate checks, 100 Darcy
tests of \texttt{list\_add}, the EL2 suite and CN verification of
\texttt{list\_add}. All stage logs are retained. The full original script subsequently completes in 22,989 seconds with eight
visible CPUs. Both trial CSVs contain all 250 prescribed observations: 220
reported detections in Darcy's symbolic mode and 30 in Bennet's random mode.
These are reproduction observations under the authors' unequal budgets
(180 and 30 seconds), with overwritten per-trial logs; they do not constitute
a matched detection comparison with \tespec{}. EL2 reports all tests successful,
and all 57 expected CN targets have individual pass markers. The verification
wrapper returns Error 123 because its worker ends with a false conditional
following a successful target, yielding exit status 1; this shell behavior
is independently reproduced. Evidence is in \texttt{artifact-full-20261011}.
The outer isolation process's resource counters are excluded from workload
performance measurements.

\subsection{Original allocator subjects: inventory probe}

We also run all 24 entries in the original pKVM allocator bug catalog with
Bennet and Darcy. Each subject uses the repository's prescribed base revision,
workaround patch and individual bug patch in an isolated worktree. This first
probe uses seed zero, ten requested samples and a 120-second limit covering
generation, compilation and execution.

\begin{center}
\small
\begin{tabular}{lrrrrrr}
\toprule
Generator & Failures & Passed & Gen. failed & Timeout & Other & Total\\
\midrule
Bennet & 3 & 0 & 21 & 0 & 0 & 24\\
Darcy & 19 & 4 & 0 & 1 & 0 & 24\\
\bottomrule
\end{tabular}

\end{center}

Failures in the table are tool-reported test failures. The original driver's
disabled replay option suppresses concrete-input printing; these observations
are distinct from independently replayed witnesses. Generation failures and
timeouts are retained separately. The small sample budget also limits input
generation's discard allowance, so these counts are not the original paper's
long-trial detection rates. All 24 Darcy controls, obtained by reversing only
the selected bug patch, pass the same requested sample budget. Pairing by case,
generator and trial yields 19 failures with passing controls, four pairs of
sampled passes and one inconclusive buggy run. Commands and seeds match within
each pair; changed predicates can nevertheless generate different concrete
inputs. A matched \tespec{} translation of these original
allocator subjects remains to be established.

We separately execute the repository's unchanged native allocator harness using
Fulminate from the April CN snapshot and Clang 18. Each run records checkout,
patching, preprocessing, instrumentation, compilation and execution as separate
stages. All seven partial-specification variants pass, as does their complete
control. Among the 24 bug-catalog variants, 13 trigger runtime check failures
and 11 pass this harness; the complete control also passes. All 33 runs reach
execution, so the failure count contains no build or instrumentation failures.
These are single-harness observations under the recorded source snapshot, with
the original embedded seed, distinct from generated-input trials and from an
exact replay of the published artifact. The archive retains all stage logs,
hashes, commands, elapsed times and peak process memory measurements.

Architecture matching is a prerequisite for translating the allocator. The
packaged executor uses a 32-bit frontend; selecting a 64-bit pointer option
alone originally left \texttt{unsigned long} at 32 bits. We added a target
consistency check and built a separate candidate using the upstream frontend's
existing 64-bit compilation option. Four size probes agree with Clang:
pointers and unsigned long occupy eight bytes, doubly linked nodes sixteen,
and explicit-padding chunk headers thirty-two. All eight generated return and
entry-satisfiability VCs have closed proofs independently accepted by
\texttt{coqchk} against an upstream parameterized logical runtime. These probes
establish selected ABI prerequisites; byte-level ownership translation and
the allocator's execution dependencies remain required for the matched port.

A separate LP64 pointer-cell test exercises ownership beyond these layout
queries. Both addresses remain symbolic; the body reads and returns the stored
pointer. The automatic run completes with an unresolved Require obligation.
An agent constructs an eight-byte mathematical heap witness and proves both
original VCs against the actual 64-bit storage definitions. Independent
certificate replay accepts both proofs with no global assumptions. The checked
pipeline takes 440.935 seconds in a run concurrent with other proof checks;
the exploratory agent construction reports 36 seconds and one compilation.
This establishes a 64-bit primitive ownership proof path while broader
allocator translation and isolated cost measurements remain pending.

A candidate field-address translation adds frontend-computed byte offsets to
symbolic bases. A 28-run matrix crosses seven layout and negative-control
cases with two target widths and two VC policies. GCC static assertions
independently validate the selected offsets; every emitted layout-marker check
passes. On 32 bits, ten positive runs prove automatically, two wrong
postconditions remain unresolved, and two overlapping-union preconditions are
rejected as unsatisfiable. Extending necessary interval constraints to 64-bit
field addresses rejects both union-overlap cases on that target as well;
the other twelve 64-bit entry checks remain unresolved automatically.
A separate two-field example has all three original VCs independently replayed
on each target, including a satisfiability proof constructing disjoint
byte cells without field-address axioms. These results distinguish successful
layout translation from completed specification tests. Imported predicates
with abstract field-address semantics require a compatible interpretation
before this candidate can become the default translation.

A separate re-elaboration experiment retains the original QCP list predicate
definitions and interprets their address notation through the exported layout.
The generic address evaluator, nested-field and array-index probes, pointer
stride probes, and a definitional unfolding check pass independent kernel replay
on both target widths. An unknown-field probe fails elaboration as intended.
The two-node list Require additionally passes the complete VC controller on
both target widths with no transitive assumptions. The input fixes logical contents and leaves
the pointer unbound. This declaration-only test establishes precondition
satisfiability; execution, safety and return checks are not run. The archived
experiment distinguishes this selected computational interpretation from
equivalence to a previously compiled predicate using abstract addresses.

The same layout generator also supports two further original predicate
families: a three-node binary tree and a two-node doubly linked queue. Their
precondition VCs pass the independent controller on both 32- and 64-bit
targets, with empty assumption sets. The tree includes two recursive branches; the queue includes head/tail
fields and forward/backward links. On both target widths, moving the fixed
logical data from Require equalities into the public JSON binding interface
produces byte-identical original Goal files and identical layouts. The tree
input uses nested constructor objects, and the queue uses a JSON list. These
are recursive-memory and instantiation checks, separate from evidence of an
expressiveness advantage over an executable monitor.

Negative counterparts equate two symbolic roots while requiring separated
copies of the same nonempty recursive predicate, using the public binding
interface for logical contents. The list, tree and queue cases each pass
independent refutation replay on both 32- and 64-bit targets with empty
assumption sets. The controller
merges these certificates into Require=\texttt{UNSAT} and
\texttt{INVALID}, preserving declaration-only scope. The contradiction arises
from duplicate ownership of the root's first byte, not from a chosen numerical
address. These checks validate negative initial-domain classification; the
six target/family combinations are independently accepted. Sources, manifests, proof
attempts and controller logs are archived separately from the positive cases.

\subsection{Independent replay of a complete heap VC group}

For the cell-increment case with logical contents $3$ and an unbound pointer,
an exploratory proof-agent run produced proofs for all four original VCs:
two arithmetic safety conditions, the return entailment, and initial
satisfiability. A separate controller pins the generated goal and dependencies,
checks the complete original module interface, audits transitive assumptions,
and runs \texttt{coqchk}. Automatic proof placeholders are excluded from this
certificate. The satisfiability proof supplies a mathematical heap witness;
the pointer remains unbound in the test itself.

All four proofs compile and pass independent replay. The return and
satisfiability lemmas are closed under the global context; the safety statements
depend on QCP's existing local-address interpretation parameter. Full-library
kernel replay takes approximately 130 seconds in this exploratory run.
An earlier 120-second replay budget produced UNKNOWN and is retained in the
data. Agent-reported proof construction took 139 seconds and ten candidate
compilations; token usage was unavailable and is recorded as missing.
These observations establish the certificate path for one generated VC group;
the controlled agent study measures its cost and coverage across subjects.

A second heap example leaves a genuine residual obligation: a function squares
an owned integer, with postcondition \texttt{is\_square} on its return and
ownership of the updated cell. An imported Rocq definition states
$\mathit{is\_square}(z)\equiv\exists k\in\mathbb{Z}.\ z=k^2$.
For the binding $n=13$, the automatic phase completes execution but leaves
the mathematical relation unresolved. The agent supplies $k=13$ and proves
all three original interface obligations, including safety and initial
satisfiability; the external controller independently replays them.
This exercises proof completion for an imported definition without supplying
a runtime decision procedure. It does not establish that an equivalent
executable checker cannot be written in a baseline system.

A paired negative case changes the postcondition to
$\mathit{is\_square}(\mathit{return}+1)$, leaving the implementation and
binding unchanged. The agent proves that $170$ is not an integer square,
refutes both branches of the original return VC, and proves the other two
interface obligations. Independent replay accepts this certificate in
132.841 seconds and records an independently checked negative proof. The refutation and initial
satisfiability proofs are closed under the global context. This result certifies
the falsity of the generated VC; a confirmed execution violation additionally
requires a reachable-state witness. The automatic UNKNOWN result is therefore
retained separately. The exploratory agent attempt reused earlier heap lemmas
and reported 30 seconds and one candidate compilation; token usage is missing.

A third case places the imported relation at intermediate boundaries. Starting
from a cell containing zero, a loop calls a helper three times to add $2i+1$;
the helper's postcondition states only that the cell contains some square.
Both execution modes retain the sequence $0,1,4,9$ despite this weak contract.
Boundary checking adds four invariant checks, three assertions and three call
contract checks. Ten intermediate entailments and the return entailment remain
unresolved automatically. An agent proves all 44 original interface obligations,
including safety, call-precondition matching and initial satisfiability, and
the independent checker accepts the group in 137.473 seconds. The ten
entailments, return and satisfiability proofs are closed under the global
context. The remaining proofs use the existing local-address interpretation
parameter. This case exercises the combination of actual calls and loops,
state-preserving boundaries and proof completion for imported predicates.
Its exploratory proof construction reused earlier lemmas and took an
agent-reported 51 seconds and two candidate compilations.

A fourth case compares an original assertion with a proved equivalent runtime
form. A loop computes the least nonnegative integer $r$ whose square is at
least the initial cell value $n$, then stores and returns $r$. Its postcondition
retains $0\le r\land\forall k\ge0.\ (r\le k\leftrightarrow n\le k^2)$.
For $n=13$ and a symbolic cell address, \tespec{} executes four iterations and
checks five invariant states. Automatic checking leaves the imported relation
unresolved; an agent proves all 26 original VCs, independently replayed in
132.063 seconds. The exploratory construction reuses earlier heap lemmas and
reports 21 seconds and one candidate compilation.

On the pinned October CN version, Fulminate instrumentation rejects the original
quantification because it cannot infer its upper bound. An identical C program
with the equivalent postcondition
$0\le r\land n\le r^2\land(r=0\lor(r-1)^2<n)$ compiles and passes native
checking on the same input. \texttt{LeastRootRewrite.v} proves this equivalence
over mathematical integers. The experiment therefore isolates direct assertion
support while retaining a successful rewritten-monitor control. CN universal
verification leaves multiplication safety unproved for both forms and reports
that nonlinear arithmetic is treated as uninterpreted; those failures are
recorded separately from instrumentation outcomes and execution violations.

\paragraph{Proof replay after entry isolation.}
The previously archived complete cell proof and square-loop boundary proof
both replay against freshly generated VCs after the entry-isolation change.
The latter completes the whole-run controller and retains checks for loop
invariants, assertions and call contracts. These are compatibility replays of
existing proofs, not additional proof-agent authoring trials.

\paragraph{Array contract translation.}
A two-cell array update provides a further paired target. Its final contracts
use typed offsets $p+0$ and $p+1$. CN verifies the valid contract, and Fulminate
and \tespec{} accept the matching concrete/logical inputs. The incorrect return
condition fails at runtime in Fulminate and remains UNKNOWN in \tespec{}.
Unmodified QCP discharges six target obligations once the native main fixture
is removed; the target's bytes are unchanged. The mutated target produces a
fatal return-obligation inconsistency diagnostic. Earlier indexing and offset
translation errors are archived separately and do not establish tool coverage
limitations. Full-unit failures after completing the target are distinguished
from target-only verification in \texttt{array-cells-20261011}.


\paragraph{Candidate and dependency regressions.}
The candidate executor with target-layout metadata passes the complete
1,322-case matrix after entry isolation, with the same status counts as the
packaged executor and zero false passes among 779 negative checks. It remains
a separately pinned candidate. A dependency replay additionally imports the
square relation through a module using \texttt{Require Export}; staging
includes both modules and excludes a fictitious import inside a nested
comment. The existing whole-instance proof passes independent replay and
run-level reconciliation. This tests transitive definition availability without
counting a new logical property or proof-agent success.

\paragraph{Complete-library layout interpretation.}
Three further compatibility replays retain the complete original singly linked
list, binary-search-tree, and doubly linked queue libraries (222, 1,696, and
223 lines). Importing the generated computational field-layout interpretation
after the original notation scope allows all three libraries, including their
existing proofs, to compile. Their two-node list, three-node tree, and two-node
queue instances then pass independent whole-run certificate checking with
empty global assumption sets. The tree proof required an explicit integer
notation scope in the candidate; its initial failed compilation is retained.
These 32-bit declaration-only cases establish Require satisfiability under the
selected interpretation. They test complete-library compatibility, rather
than adding independent proof examples or establishing equivalence to the
legacy abstract-address semantics. The archive
\texttt{proof-audit/full-layout-libraries} contains the original and transformed
sources, frozen goals, candidates, reports, and compressed stage logs.

\paragraph{Recursive traversal of a linked heap.}
A further matched target recursively computes the length of a singly linked
list while preserving its contents. CN's recursive content predicate and QCP's
original list predicate describe the same two-node fixture, with an upper
length bound of 100. CN verifies the correct contract after an explicit
unfolding hint for its mathematical length function; Fulminate accepts the
matching execution. Changing the specified result to length plus one produces
an unprovable CN constraint and a Fulminate runtime failure. In both execution
modes, \tespec{} completes the recursive traversal and automatically checks
safety and the correct return condition; boundary mode also checks the actual
recursive call contracts. Initial satisfiability remains UNKNOWN in the
original-library run, so neither valid run is counted as an overall success.
Unmodified QCP produces four automatic and four residual target obligations
for each variant. A target-only control separates these results from a later
failure in the native main fixture. The archive \texttt{list-length-20261011}
retains final runs and earlier translation corrections; they represent one
subject, and their concurrent timings are diagnostic measurements.

The linked-list regression additionally varies length from zero through five
and changes the recursive increment from one to two. All 24 symbolic runs
(two bodies, six inputs, two modes) finish. Correct-body postconditions pass
throughout; the mutant passes only on the empty list and leaves a residual
postcondition for every nonempty input. Twelve separately compiled native
fixtures agree with the arithmetic oracles $n$ and $2n$. Initial satisfiability
remains unresolved in these original-library runs, so the grid contributes
phase-level regression evidence rather than overall success or confirmed
symbolic counterexamples. Inputs, native fixtures, and reports are archived in
\texttt{list-length-grid-20261011}.

\paragraph{Proof completion for the recursive linked heap.}
A separate interpretation experiment re-elaborates the complete original list
library using explicit frontend field offsets and checks the actual two-node
recursive execution in boundary mode. The proof agent closes all 19 original
interface obligations, including initial satisfiability, safety, recursive
call boundaries, and final return. Independent controller compilation,
assumption auditing, and kernel replay accept the complete group and merge
all execution phases to success. The initial satisfiability theorem is closed;
the complete group additionally depends on the runtime's local-variable
address parameter and functional and propositional extensionality, with no
unexpected assumptions. Authoring takes 236 seconds over seven compilation
attempts; independent checking takes 140.12 seconds. These concurrent
development measurements include two terminated tactic attempts, and token
telemetry is unavailable. The evidence is archived in
\texttt{proof-audit/list-execution}. This result establishes execution under
the selected computational layout; the original-library automatic outcome
remains separately recorded as UNKNOWN.

\paragraph{Binary-tree execution barrier.}
An additional coverage probe recursively counts both subtrees of a three-node
binary tree. Fulminate accepts the valid fixture and detects the mutated return
condition after renaming ghost constructors to avoid a generated-C identifier
collision. The first CN annotation leaves the nonnegative subtree-size fact unproved
at the recursive call. Adding nonnegativity assertions for both recursive
predicate outputs lets CN verify the valid contract and reject the mutated
postcondition as unproved; the same fixtures pass and fail Fulminate runtime
checks respectively. Structural nonnegativity is separately proved in Rocq.
This retains the repaired baseline result alongside the initial annotation failure. \tespec{} fails at the first child-field read because the
original external \texttt{store\_tree} predicate remains opaque in its
execution state. Providing the full Rocq definition as a proof dependency
does not automatically expose it to executable predicate normalization.
This is an execution-coverage gap, distinct from the linked-list example's
initial satisfiability obligation. Both annotation attempts and the exact
failure states are retained in \texttt{tree-size-20261011}; the probe is
reported separately from the reviewed paired table.

\paragraph{Definition-driven unfolding probe.}
A subsequent experiment translates constructor clauses of imported heap
predicates into QCP strategy proposals and Rocq branch equations from a shared
intermediate representation. The same translator handles the original tree,
list, and list-segment predicates. All six branch equations compile and pass
an independent \texttt{coqchk} run. A renamed predicate with different
constructor and field names also compiles; deleting one owned field from its
proposed expansion causes the equation proof to fail. With the generated tree
rules supplied explicitly, both execution modes complete all six recursive
calls of the three-node fixture and retain return value 3. The valid and
mutated postconditions both remain UNKNOWN: pure logical-function evaluation,
initial satisfiability, and internal contract binding still require further
work. Thus this probe removes the observed field-read barrier without
establishing a complete test verdict. The equations certify the logical branch
reconstruction; the translation to QCP strategy syntax remains part of the
trusted implementation. The mechanism is experimental and not enabled by
default. Sources, proofs, the rejected mutation, and all four execution runs
are retained in \texttt{proof-audit/rocq-predicate-lowering}.

\paragraph{Reconstructing callee model arguments.}
Further probes isolate the contract-binding gap after expansion: a callee
requires an abstract tree value, whereas the caller exposes its fields.
Goal-side constructor inference restores this value without changing the
execution state. Both literal and symbolic forms of base-case guards must be
handled. An initial 28-run grid exposes four remaining binding failures;
after supporting both forms, all 14 repeated no-expand runs infer their
callee arguments successfully. These runs cover seven tree shapes, including
empty, left- and right-skewed, and balanced trees of up to seven nodes, with
both the original traversal and a mutation that counts each node twice.
The top-level return literals in the generated VCs agree with $n$ and $2n$,
respectively. All runs remain UNKNOWN because the logical size function and
initial satisfiability still require proof; successful argument inference
is not counted as a proved contract or detected mutation. Six generated
constructor-substitution lemmas compile alongside the branch equations; all
three generated modules also pass an independent \texttt{coqchk} run.
The attempts and final grid are retained in
\texttt{proof-audit/rocq-predicate-inference}.

\paragraph{Complete binary-tree execution certificate.}
A subsequent three-node fixture combines generic constructor unfolding and
model inference with the explicit 32-bit computational layout. Actual execution
performs six recursive calls and returns 3. The independent controller accepts
all 44 original obligations, including initial satisfiability, safety, call
returns, the final return and auxiliary entailments. All return proofs retain
the complete heap branch. The assumptions are the existing address-interpretation
parameter, functional extensionality and propositional extensionality; legacy
field-address axioms are absent. Replay takes 386.405 seconds, including
261.559 seconds for \texttt{coqchk}, in this exploratory run. Evidence is in
\texttt{proof-audit/tree-execution}. The result uses the selected computational
interpretation; the earlier original-library coverage outcomes remain recorded
separately. Agent token and full attempt telemetry are unavailable.

\paragraph{Ground logical evaluation.}
A generic Rocq evaluation service computes the logical tree-size function on
the seven instantiated shapes and generates seven equality certificates.
All certificates compile with no global assumptions and pass independent
\texttt{coqchk}. A separate recursive weighted-tree function exercises negative
integer results; opaque terms remain UNKNOWN, and a corrupted computed result
fails proof checking. Connecting the certificates to the actual three-node
return VC exposes a namespace problem: an imported heap notation shadows
integer addition in the emitted goal. Restoring the integer scope in the
executor's generated headers permits the unedited, regenerated goal to compile.
Rewriting with the computed equalities then proves its return-arithmetic
subgoal with no global assumptions. This closes one real VC component;
initial satisfiability and remaining spatial obligations still determine the
overall UNKNOWN result. The candidate binary, failed attempts, successful
proofs, and per-process resource measurements are identified in
\texttt{proof-audit/rocq-ground-evaluation}. The packaged executor and the
earlier paired observations remain unchanged.

\paragraph{Audited paired observation table.}
Table~\ref{tab:paired-observations} consolidates 11 paired subjects and 102
individual tool invocations from the final development probes. The deterministic
summary validates source and log hashes and reads \tespec{} phase results
directly, independently of the smoke runner's expected-outcome flags. Repeated
CN/Fulminate invocations across \tespec{} modes remain repeated observations
of the same subject. Earlier setup attempts, deliberately invalid helper
contracts, and proof-completed replays are separate strata.
\begin{table}[t]
\centering
\small
\begin{tabular}{llll}
\toprule
Subject & CN valid / mutated & Fulminate valid / mutated & TeSpec valid / mutated \\
\midrule
alias\_shared & verified / unproved & pass / violation & pass / unknown \\
array\_cells & verified / unproved & pass / violation & pass / unknown \\
heap\_loop & verified / unproved & pass / violation & pass / unknown \\
list\_length & verified / unproved & pass / violation & unknown / unknown \\
nested\_call & verified / unproved & pass / violation & pass / unknown \\
pointer\_cell & verified / unproved & pass / violation & pass / unknown \\
recursive\_heap & verified / unproved & pass / violation & pass / unknown \\
recursive\_heap\_weak & unproved / unproved & pass / violation & pass / unknown \\
single\_cell & verified / unproved & pass / violation & pass / unknown \\
struct\_fields & verified / unproved & pass / violation & pass / unknown \\
swap & verified / unproved & pass / violation & pass / unknown \\
\bottomrule
\end{tabular}
\caption{Final paired development probes. Modes agree on these overall results; the archive retains their different annotation policies and phase verdicts. CN verifies the full precondition domain, Fulminate runs a concrete fixture, and TeSpec checks a bound symbolic instance. Unknown/unproved results are not confirmed violations. These observations exclude subsequent proof completion and are not performance trials.}
\label{tab:paired-observations}
\end{table}


\subsection{Comparative experiment design}

The main experiment will combine original Fulminate/Bennet/PLDI artifact
subjects, the complete QCP inventory, and paired common-fragment programs.
Expression-level and ownership-level contract translations require review for
semantic equivalence. Proof-oriented separating examples form a distinct
stratum so logical expressiveness is measured independently of supported C
features. Input generation and checking a supplied input are measured separately.

Four paired experiments connect each comparative claim to its mechanism.
First, original quantified or relational assertions are paired with executable
rewrites whose equivalence is checked, measuring direct assertion support and
rewrite effort. Second, mutations of assertions, invariants, and call contracts
preserve the final output while changing an intermediate logical obligation,
measuring observation coverage for the same rich property. Third, fixed logical
data with symbolic addresses are paired with concrete layouts and aliasing
controls, measuring the input domain established by a completed instance.
Fourth, normalization, unfolding, and proof completion are ablated on the same
subjects, separating automatic simplification from residual proof work.
Subjects and feature categories are registered before outcomes are collected.

The primary unit is a complete instance: initial nonemptiness, execution safety,
and all selected boundary obligations must be checked. Individual discharged
VCs measure automation within that unit. The current interface reports FAIL when an admitted original path VC is
INVALID; unresolved proof obligations remain UNKNOWN. Archived negative proof
replays retain their recorded evidence scope. Common-fragment controls quantify overhead,
and capability-separating subjects quantify the additional original assertions
that obtain complete verdicts. Together these strata measure the benefit and
cost of retaining the verification logic as the oracle.

Outcomes include certified success, confirmed violation, invalid input,
unsupported behavior, error and timeout. Ablations compare normalization,
SMT-only checking, checked ITP completion and intermediate-boundary policies.
Agent experiments record attempts, models, prompts, tokens, time, proof replay
and assumptions. Deterministic timing uses repeated isolated runs; stochastic
conditions use recorded independent trials and fixed budgets. All exclusions
and unsuccessful runs remain in the released data.

\section{Mechanized Semantic Core}

The accompanying \texttt{InstanceChecking.v} defines typed partial bindings,
their instance domains, finite symbolic execution trees, and proof obligations.
Its central theorem states that a closed tree establishes safe, terminating
behavior for every state represented at its root. Step nodes require progress
and coverage of every concrete successor; terminal nodes require safety and
the postcondition. Normalization requires semantic equivalence. Boundary nodes
require the annotation and retain the same successor assertion. An unfinished
node has no certificate.

An acceptance theorem combines this coverage result with a separate nonemptiness
proof, yielding nonvacuous instance correctness. Further theorems establish
domain inclusion under stronger bindings and preservation of joint assertion/
return observations. A mechanized counterexample shows why existential
projection that preserves the closed assertion can still corrupt return-value
checking. The development compiles with Coq 8.20.1; its principal theorems have
no additional global axioms and its compiled library is independently replayed
with \texttt{coqchk}.

The negative-result model requires a bad state reachable from a member of the
initial instance domain. It proves that an accepted execution certificate
excludes such a witness, and that a bad model of a boundary assertion yields
a counterexample when every model of that assertion is reachable. A separate
mechanized example demonstrates why refuting an implication over an
overapproximation alone is insufficient: the implication may fail only on
states that no execution reaches.

The initial oracle has a separate four-statement mechanization in
\texttt{PreconditionOracle.v}. For the emitted proposition
$\exists\eta.\,\neg\forall h.\,(R(\eta,h)\rightarrow\bot)$, a concrete logical
witness implies the VC constructively; the converse uses excluded middle as
an explicit premise. Negating the VC is constructively equivalent to emptiness
of the entire initial domain. A counterexample preserves a satisfiable
alternative when another branch is refuted. A further eight statements in
\texttt{AffineOffsets.v} establish common-base integer cancellation and
byte-interval simplification, including preservation of other query constraints
and rejection of overlapping offsets. The observation module adds nine
statements for exact local framing, visible-heap uniqueness, resource retention
and rejection of an unmentioned user cell, including witness-dependent frames. A further five statements establish finite integer-list equality reflection,
append/reverse comparison and preservation of owned heap observations.
Seven further statements establish typed-witness inhabitance, empty-type
refutation, self-equality preservation and safe elimination for integer/list
types. The budgeted definition-normalization module adds seven statements preserving
joint observations, nonemptiness, validity and countermodels under licensed
rewrites. Across the ten standard-library modules, 82 principal statements have closed global contexts and pass
independent kernel replay. The classical premise is part of the stated
positive equivalence, not an additional global assumption.

These theorems parameterize the operational relation and assertion
interpretation. Establishing that the C executor generates obligations satisfying
those interfaces remains an implementation-refinement obligation, including
call-frame lifetime, C arithmetic and layout, substitution and SMT encoding.
The generic theorem and the implementation tests are separate evidence.

An additional nested-field regression executes two function calls and observes
three invariant states and a final assertion without replacing the execution
state. Its no-expand controller checks the definition equations and their
association with the program's memory interpretation before execution.
Local scalar framing and normalized field-strategy matching preserve the
existing heap tests; the current packaged executor passes 274 tests, with ten additional
tests skipped. Wrong field values, wrong local values and omitted user resources
remain unresolved rather than becoming false successes. Local-pointer alias observations now select current goal witnesses before
framing, including double indirection, rebindings and symbolic branches. These results measure implementation
coverage; they establish no comparative advantage without a matched baseline.

A further dependency-context check unfolds the original functional-queue
predicate under its original module identity. Its five branch and standard-list
identity obligations pass independent replay without manually importing the
auxiliary list module. Four retained association checks accept the matching
environment and reject constructor, relation and jointly shadowed predicate
bindings. In the last case both ordinary aliases agree, but the additional
qualified-definition identity fails. These checks validate definition selection;
they do not constitute a completed queue execution. Evidence is retained in
\texttt{proof-audit/qualified-normalization}.

Typed nested-pattern translation extends the same library scan to 23 parsed
and five unsupported declarations. Three original tree-context predicates,
including the parent-pointer variant, pass all 24 unfolding, list-constructor
identity and model-choice obligations under their original module identities.
A constructed nested-model memory read completes the controller with initial
SAT and execution safety and return PASS. Ten direct executor regressions in
the two boundary policies include actual loops and calls: six pass, two reject
inconsistent inputs, and two retain incorrect return properties as unresolved.
These outcomes exercise generic constructor translation and preserve the
distinction between a rejected input and an unclosed execution proof.
Evidence and unsuccessful preparation attempts are retained in
\texttt{proof-audit/typed-constructor-patterns}.

Nominal-witness and pure-relation support advances the same scan to 26 parsed
declarations and two layout barriers. The original \texttt{Bst.store\_map}
definitions in two BST libraries each pass an independently checked unfolding
equation, preserving existential tree witnesses and their \texttt{SearchTree}
and \texttt{Abs} relations. An empty nominal type exposed a frontend erasure
bug: its unused witness disappeared and the old executor accepted an impossible
Require. The repaired executor retains the typed witness and reports UNKNOWN
until checked proof completes the result. Two complete three-VC groups yield
INVALID for the empty type and PASS for a constructor-inhabited type.
Renaming, strategy-generated witnesses, both boundary policies and typed
assertion/invariant regressions check the preservation mechanism; the latter
continue the actual loop and prove its return despite unresolved annotations.
A fresh normalization controller also passes definition equations and program
association before the nested-model heap read completes with automatic SAT
and safety/return PASS. The execution verdict uses the ordinary solver; its
VCs have no separate kernel replay in this controller run. Evidence is retained
in \texttt{proof-audit/nominal-existentials}.

The address-term extension also parses the original pointer-segment predicate
\texttt{sllbseg}, including its recursive next-field argument. Both branch
equations pass independent replay under the original runtime module identity.
Program association still rejects its abstract layout for computational
execution. A separate explicit-layout recursive predicate completes the
no-expand controller with initial SAT, an actual heap read, safety and return
PASS. The untouched three-element insertion-sort instance executes but retains
unproved initial and relational obligations. The library scan has 18 parsed
declarations and ten unsupported declarations; neither parsing nor an unfolding
equation counts as a completed program test. Sources and all results are retained
in \texttt{proof-audit/recursive-address-arguments}.

The layout bridge contributes eighteen of these closed theorems over explicit
object and field layouts. Addresses are computed by adding byte offsets;
nested paths compose these additions. The theorems establish bounds,
containment, inherited alignment, positive addresses and byte-domain
disjointness under explicit layout conditions. In particular, disjointness of
enclosing fields implies disjointness of their contained paths, independently
of structure names or nesting depth. Connecting frontend metadata and the
generated heap predicates to this interpretation remains necessary for
certifying the frontend translation as a whole. The instance certificates
above separately establish the selected generated heap obligations.

A runtime-dependent development adds nine closed lemmas for the partial address
evaluator, including mutual correctness results for lvalue and rvalue address
expressions relative to an inductive interpretation. These lemmas are replayed
separately on both memory runtimes, alongside four layout probes and one
predicate-unfolding lemma per target. Their scope is address computation for a
given layout environment. Frontend layout correctness and C memory safety
remain distinct obligations.

\paragraph{Paired quantified intermediate observation.}
A further pair keeps the least-root routine, data binding $n=13$, and final
postcondition unchanged, and changes only the assertion before the heap write
from $\mathit{least\_root}(n,r)$ to
$\mathit{least\_root}(n,r-1)$. Both tests execute four loop iterations and
visit five invariant boundaries. Automatic execution establishes initial SAT
and safety, leaving the quantified observation and postcondition unresolved.
The fixed controller checks all 27 original obligations of the positive case
and produces an overall PASS report. For the negative case, it checks the
negation of the complete assertion VC, including both generated alternatives,
and proves the other 26 obligations. This archived replay records a negative VC proof, rather than
an execution FAIL in the archived controller: a reachable execution counterexample certificate was still
required. Full certificate checks take approximately 124.45 and 125.86 seconds,
respectively; the positive end-to-end report records its separate controller
replay. Proofs in this audit were supplied during development, so these costs
do not measure automatic proof-agent performance. CN translations of the
intermediate assertions remain unresolved and supply no comparative advantage
result. Sources, diagnostics, checked proofs, and final reports are archived in
\texttt{proof-audit/intermediate-quantifier-pair}.

\paragraph{Checked pure consequences during recursive unfolding.}
The original imported \texttt{sll\_first\_value} instance binds the head to
$1$ and the logical tail to $[2,3]$. Its numerical pointer address remains
unbound in a no-expand replay. The imported left-unfolding strategy can discard
the non-null conditions of tail nodes. We certify witness-independent pure
consequences against the original definition, retaining each predicate, and
apply these consequences between strategy steps as further applications
appear. Independent equation and program-association checks pass. The actual
read, safety and return obligation pass automatically; initial nonemptiness
remains unresolved, so the overall result is \textsf{UNKNOWN}. A companion
expand replay with a bound head address also passes its return obligation.
The symbolic-address controller takes 130.28 seconds including cold kernel
replay; automatic execution takes 0.0145 seconds with 5,348 KiB peak process
RSS. These are development measurements, not proof-agent performance results.
The packaged regression replay passes 301 tests, with 12 environment-dependent
skips. Library inventories identify eligible pure-fact proposals in 14 of 28
human-library definitions and 16 of 38 LLM-library definitions; these counts
measure preparation, not successful proofs or program coverage.
Evidence is archived under \texttt{proof-audit/pure-fact-projection}.

\paragraph{Measured benefit of optional external normalization.}
Matched ordinary-strategy baselines for the two configurations above leave
both initial existence and the root return obligation unresolved. Checked
pure consequences remove the root return obligation in each configuration,
reducing root residual obligations from two to one; split subgoals are not
counted separately. Both overall outcomes remain \textsf{UNKNOWN}, yielding
zero additional complete passes. These are two configurations of one program
and one logical instance, not two independent benchmarks. Baseline executor
times are 0.0120 and 0.0125 seconds; projection executor times are 0.0105 and
0.0145 seconds. Single runs do not establish a speedup. Cold checked-projection
controller costs are 129.41 and 130.28 seconds. A full-expansion attempt spends
123.99 seconds before the computational field-interpretation gate rejects it;
it does not execute and demonstrates no solver benefit. Accordingly, imported
definitions remain opaque by default. The results support a scoped pure-fact
optimization for this fixture, rather than mandatory full expansion or a
claim of dataset-wide improvement. Pairing checks and reports are archived in
\texttt{proof-audit/external-expansion-benefit}; historical projection reports
lack the baseline report's full provenance block, as recorded in that audit.

\paragraph{Default workflow on original external programs.}
A fresh replay without added external definition rules covers eight instances:
empty and three-element original insertion-sort inputs, a nonempty original
queue-dequeue input, and a declaration-only reverse input, each in both modes.
All four sort executions complete actual loops and helper bodies, but their
initial existence and functional obligations remain unresolved. Both reverse
instances produce only initial-existence tasks and report implementation
execution and postcondition checking as not run. Both queue instances report
an error on reading a field hidden by the opaque queue predicate; the imported
queue strategies have no left exposure rule. Thus these runs yield six unknown
and two error results, with no complete pass. They distinguish executed but
unproved tests from absent bodies and inaccessible resources. ITP preparation
also exposed a controller omission: QCP emits a spatial split child as well as
numbered pure children. The repaired reconciler recognizes both forms only
when the complete original parent belongs to the certificate interface.
Prepared sort groups retain all original obligations: 5 and 18 in expand mode,
6 and 42 in no-expand mode for empty and three-element inputs respectively.
Preparation is not proof completion. Sources, reports, initial failed
preparations and corrected tasks are in
\texttt{proof-audit/default-external-workflow}.

\paragraph{Broad original-corpus smoke coverage.}
A default-strategy replay inventories all 81 sources in the original human
and LLM demo directories and selects 327 specifications. Ninety-three require
explicit typed bindings; the remaining specifications generate 1,404 smoke
cases using synthetic sizes zero, one and two, both modes and symbolic
pointers. The fixed executor reports 448 passes, 726 unknowns, 128 invalid
inputs and 102 errors. Of the passes, 212 are execution cases and 236 are
specification-only; the latter do not verify implementations. The unknown
results include 61 timeouts, 55 loop-budget outcomes and 208 missing-body
outcomes. These generated inputs have no independent outcome oracle and are
engineering coverage, not a comparative benchmark. A generic repair prompted
by six named-invariant routing errors lets all six affected original
bubble-sort cases finish execution; their final outcomes remain unknown.
The complete fixed corpus denominator, generated artifacts and subsequent
repair replays are archived in \texttt{proof-audit/default-corpus-20261012}
and \texttt{proof-audit/unrouted-observation} respectively.

\paragraph{Eliminating spurious pointer-difference paths.}
Matched original array-sum inputs $[1]$ and $[1,2]$ leave addresses symbolic.
The previous executor fails both modes by retaining a guard whose numerator
subtracts two offsets of the same base; the resulting spurious path reads
past the array. Generic shared-base difference normalization removes this
path. Standard \texttt{sublist} evaluation then reduces closed slices in the
visited invariants, allowing QCP's existing sum rules to close them. Both
inputs pass completely in both modes. No added external spatial-definition
scope or program-specific sum rule is used. The generic regressions preserve
unknown list spines, symbolic elements, negative-bound semantics and distinct
pointer bases. Two additional closed arithmetic statements establish
common-base difference and exact typed quotient cancellation. Six examples
against QCP's actual imported slicing definition pass independent kernel
replay. These are component proofs, not a whole-executor C refinement.
Evidence and the matched frozen old/final reports are in
\texttt{proof-audit/shared-base-and-sublist}. The intermediate build was not
frozen and is not reported as a reproducible ablation arm.

\paragraph{Separating inconsistent inputs from memory-read limits.}
A restored-executor trace for an original AVL rotation shows that native rules
already expose node fields. Its synthetic zero binding requires
$store\_non\_empty\_tree(0)$; the original definition contains $0\ne0$.
Two standalone Rocq lemmas refute this predicate and the complete instantiated
rotation Require for every symbolic root, followed by independent kernel replay.
Their assumptions are the original abstract address evaluator and functional
extensionality. This explains that one input without establishing outcomes for
other AVL bindings. The automatic report remains initial UNKNOWN and execution
ERROR: attaching the standalone refutation to an immutable generated Require
VC remains necessary. Failure before VC emission is a pipeline barrier, not
proof that the original Require is satisfiable.
